\section{Conclusion}

This thesis draws from software design methodologies to develop, structure, and verify security models. This research demonstrates the value of modularity and software-inspired techniques for building cybersecurity models. 

Modularity provided through compositional game theory gives clear advantages for conceptually adapting or composing models together. The addition of a structured development process makes these advantages readily available for practitioners. As an "applied-applied" category theory project, we showed that using the category of open games and its Haskell implementation can serve effectively as a backend for a platform that allows cybersecurity modellers to experiment, collaborate, and save their models. The introduction of design patterns to compartmentalize certain operations makes the open games software more accessible and straightforward-to-use. Then, establishing an iterative, test-driven process to continuously update the resultant models then leverages these design patterns to organize a repository of models.  

Alongside active development and improvement of the open-game-engine, we hope to continue to expand the capabilities of this workflow. More active pipelines and a behavioural study of users of this workflow could propel interest in both cybersecurity and microeconomic game-theoretic modelling, but also a wider adoption of the open-games-engine as a viable modelling tool in other disciplines. Further development of other tools to promote greater synergy between model conceptualization and desired specifications is also shown through the CyberVITBridge. 